\section{Verifica\c{c}\~ao}
\label{sec:verificacao}

Cada pe\c{c}a foi verificada contra um or\'aculo \emph{exato} antes de a
pr\'oxima ser constru\'ida. Os arreios est\~ao em
\texttt{higflow/tests/front-tracking} e devolvem c\'odigo de sa\'ida: s\~ao
port\~oes, n\~ao impress\~oes.

\subsection{B1: advec\c{c}\~ao cinem\'atica reversivel}
\label{sec:b1}

O teste de Rider--Kothe~\cite{rider1998}: um c\'irculo advectado no campo de
v\'ortice \'unico
\begin{align}
  u &= \phantom{-}\sin^{2}(\pi x)\sin(2\pi y)\cos(\pi t/T),\\
  v &= -\sin(2\pi x)\sin^{2}(\pi y)\cos(\pi t/T),
\end{align}
que \'e livre de diverg\^encia e se inverte em $t=T/2$, de modo que a frente deve
voltar \`a forma inicial em $t=T$. Dois or\'aculos: a \emph{conserva\c{c}\~ao} de
\'area ao longo de toda a corrida, e a \emph{revers\~ao} em $t=T$.

\pgfplotstabletypeset[
  columns/nmarc/.style={column name={$N$}, int detect},
  columns/dt/.style={column name={$\Delta t$}, fixed, precision=3},
  columns/conssem/.style={column name={cons.\ sem cir.}, sci, sci zerofill, precision=2},
  columns/revsem/.style={column name={rev.\ sem cir.}, sci, sci zerofill, precision=2},
  columns/conscom/.style={column name={cons.\ com cir.}, sci, sci zerofill, precision=2},
  columns/revcom/.style={column name={rev.\ com cir.}, sci, sci zerofill, precision=2},
  columns/nfinal/.style={column name={$N$ final}, int detect},
  every head row/.style={before row=\toprule, after row=\midrule},
  every last row/.style={after row=\bottomrule},
]{\dat{b1-rider-kothe.dat}}

A advec\c{c}\~ao \'e de segunda ordem: sem cirurgia, a revers\~ao cai
\num{3.2e-4} $\to$ \num{4.0e-5} $\to$ \num{5.0e-6}. Com cirurgia, a
conserva\c{c}\~ao de \'area fica \emph{uma ordem de grandeza melhor} sob
esticamento, porque manter $\Delta s \approx h$ preserva a resolu\c{c}\~ao do
filamento; em troca, a revers\~ao perfeita se perde, o que \'e o compromisso
correto.


\begin{figure}[htbp]
\centering
\begin{tikzpicture}
\begin{axis}[
  width=11cm, height=8cm, axis equal image,
  xlabel={$x$}, ylabel={$y$}, grid=major,
  xmin=0.1, xmax=0.9, ymin=0.3, ymax=1.0,
  legend pos=south west, legend cell align=left, legend style={font=\small},
]
\addplot[thick]              table[x=x,y=y] {\dat{formas/b1_0.dat}};
\addlegendentry{$t=0$ (c\'irculo)}
\addplot[thin, dashed]       table[x=x,y=y] {\dat{formas/b1_2.dat}};
\addlegendentry{$t=T/2$ (filamento)}
\addplot[thick, densely dotted] table[x=x,y=y] {\dat{formas/b1_4.dat}};
\addlegendentry{$t=T$ (retorno)}
\end{axis}
\end{tikzpicture}
\caption{Teste revers\'ivel de Rider--Kothe. A frente estica em filamento at\'e
$t=T/2$, quando o campo se inverte, e retorna \`a forma inicial em $t=T$. A
cirurgia acompanha o esticamento inserindo marcadores (256 $\to$ 680) e os
remove no retorno (654). A \'area fechada varia \num{6e-5} ao longo de todo o
percurso.}
\label{fig:b1}
\end{figure}

\paragraph{O que custou uma itera\c{c}\~ao.} A primeira vers\~ao da cirurgia
removia marcador de segmento curto \emph{sem guarda de curvatura}, e cortava
canto: o erro de \'area passou a ser dominado pela pr\'opria cirurgia, com a
assinatura inconfund\'ivel de o erro m\'aximo coincidir com o erro final. A
guarda descrita na Se\c{c}\~ao~2.3 corrigiu, e o teste mostrou a melhora.

\subsection{Curvatura}

Contra dois or\'aculos anal\'iticos: c\'irculo ($\kappa=1/R$ em todo marcador) e
elipse ($\kappa(t) = ab/(a^{2}\sin^{2}t + b^{2}\cos^{2}t)^{3/2}$).

\pgfplotstabletypeset[
  columns/n/.style={column name={$N$ (c\'irculo)}, int detect},
  columns/circulo/.style={column name={erro rel.\ c\'irculo}, sci, sci zerofill, precision=2},
  columns/nelipse/.style={column name={$N$ (elipse)}, int detect},
  columns/elipse/.style={column name={erro rel.\ elipse}, sci, sci zerofill, precision=2},
  every head row/.style={before row=\toprule, after row=\midrule},
  every last row/.style={after row=\bottomrule},
]{\dat{curvatura.dat}}

No c\'irculo a curvatura \'e \textbf{exata} em precis\~ao de m\'aquina --- tr\^es
pontos sobre um c\'irculo definem esse mesmo c\'irculo. Na elipse converge em
segunda ordem. Esse resultado \'e o que autoriza, mais adiante, usar
$\sigma/R$ como refer\^encia exata para o salto de press\~ao: a curvatura que a
for\c{c}a \eqref{eq:tensao} aplica numa gota circular \'e $1/R$ sem erro de
discretiza\c{c}\~ao.

\subsection{Espalhamento da for\c{c}a}

Dois or\'aculos sobre uma gota em repouso, ambos em precis\~ao de m\'aquina:
\emph{conserva\c{c}\~ao} --- $\sum_{\text{malha}} \mathbf{F}\,h^{2}$ deve igualar
$\sigma\sum_k \kappa\mathbf{n}_k\,\mathrm{d}s_k$, porque o n\'ucleo de Roma soma
1 (parti\c{c}\~ao da unidade) --- e \emph{equil\'ibrio}, $\oint \kappa\mathbf{n}\,
\mathrm{d}s \to 0$ para curva fechada, que \'e a for\c{c}a l\'iquida nula exigida
pela lei de Laplace.

\pgfplotstabletypeset[
  columns/nmarc/.style={column name={$N$}, int detect},
  columns/conservx/.style={column name={conserv.\ $x$}, sci, sci zerofill, precision=2},
  columns/conservy/.style={column name={conserv.\ $y$}, sci, sci zerofill, precision=2},
  columns/equilib/.style={column name={equil\'ibrio $\lvert\mathbf{F}\rvert$}, sci, sci zerofill, precision=2},
  every head row/.style={before row=\toprule, after row=\midrule},
  every last row/.style={after row=\bottomrule},
]{\dat{tensao-conservacao.dat}}

\paragraph{Honestidade sobre o segundo or\'aculo.} O equil\'ibrio dar
m\'aquina-zero vem da \emph{simetria do c\'irculo}; uma curva fechada
assim\'etrica daria um valor finito convergindo. Como o c\'irculo \'e exatamente
o caso de Laplace, m\'aquina-zero \'e a resposta certa aqui --- mas o teste forte
desta subse\c{c}\~ao \'e a conserva\c{c}\~ao, que n\~ao depende de simetria.

\subsection{Advec\c{c}\~ao acoplada}

Com o solver no circuito, a frente passa a se mover com a velocidade
interpolada da malha. A gota est\'atica serve de or\'aculo: ela deve ficar parada,
circular e com \'area preservada, j\'a que o campo projetado \'e livre de
diverg\^encia.

\pgfplotstabletypeset[
  columns/passo/.style={column name={passo}, int detect},
  columns/n/.style={column name={$N$}, int detect},
  columns/area/.style={column name={\'area}, fixed, precision=8},
  columns/driftrel/.style={column name={$\Delta A/A$}, sci, sci zerofill, precision=2},
  columns/deriva/.style={column name={deriva centr.}, sci, sci zerofill, precision=2},
  columns/circ/.style={column name={circularidade}, fixed, precision=6},
  columns/unidade/.style={column name={part.\ unidade}, sci, sci zerofill, precision=2},
  columns/umarc/.style={column name={$\max\lvert u_{\text{marc}}\rvert$}, sci, sci zerofill, precision=2},
  every head row/.style={before row=\toprule, after row=\midrule},
  every last row/.style={after row=\bottomrule},
]{\dat{adveccao-acoplada.dat}}

A \emph{parti\c{c}\~ao da unidade} em \num{6.9e-15} \'e o or\'aculo forte: os
pesos somam 1 em precis\~ao de m\'aquina, o que prova que um campo uniforme seria
interpolado exatamente e pega suporte incompleto --- que daria velocidade pequena
demais, em sil\^encio.

\paragraph{A checagem que evitou um falso verde.} Uma gota que fica parada \'e
\emph{exatamente} o que se veria se a interpola\c{c}\~ao devolvesse zero em
sil\^encio: \'area perfeita, deriva nula, circularidade constante. Da\'i a
\'ultima coluna. Ela vale \num{2.9e-4} --- n\~ao nula --- e vale
\emph{exatamente zero no passo 0}, quando a condi\c{c}\~ao inicial \'e
$\mathbf{u}=0$. O n\'umero acompanha o campo, logo a advec\c{c}\~ao est\'a de
fato operando.

\paragraph{O que estes testes n\~ao exercitam.} A gota est\'atica n\~ao deforma:
$N$ permanece constante e \textbf{a cirurgia nunca dispara} no regime acoplado.
O caveat de $\kappa=0$ em marcador rec\'em-inserido segue, portanto, sem
verifica\c{c}\~ao com o solver no circuito. Ambos entram em cena em B3 e B4.
